\section{Bukti Langsung dan Kontraposisi}

\textbf{Bukti langsung} adalah metode paling alami untuk membuktikan implikasi $p \rightarrow q$ \cite{rosen2019,forallx}. Prosedurnya: asumsikan bahwa $p$ benar, lalu dengan langkah-langkah logis yang valid tunjukkan bahwa $q$ juga harus benar. Karena implikasi $p \rightarrow q$ hanya salah ketika $p$ benar dan $q$ salah, membuktikan bahwa $q$ benar apabila $p$ benar cukup untuk menetapkan kebenaran implikasi tersebut. Pendekatan ini cocok ketika kita memiliki cara yang relatif langsung untuk menurunkan $q$ dari $p$, misalnya dengan menggunakan definisi, aksioma, atau teorema yang telah terbukti \cite{levin_discrete,copi2014}. Gambar~\ref{fig:direct-vs-contrapositive} membandingkan alur kedua metode.

\begin{figure}[htbp]
\centering
\begin{tikzpicture}[node distance=0.9cm, font=\small, box/.style={rectangle, draw=black, rounded corners=3pt, minimum width=2.4cm, align=center, fill=#1}]
  % Bukti langsung
  \node[box=blue!12] (d1) at (0,2) {Buktikan $p \to q$};
  \node[box=blue!10] (d2) at (0,1) {Asumsikan $p$ benar};
  \node[box=blue!10] (d3) at (0,0) {Turunkan $q$};
  \draw[thick, ->, >=stealth] (d1) -- (d2);
  \draw[thick, ->, >=stealth] (d2) -- (d3);
  \node[below=0.15cm of d3] {\textbf{Bukti langsung}};
  % Kontraposisi
  \node[box=green!12] (c1) at (4.5,2) {Buktikan $p \to q$};
  \node[box=green!10] (c2) at (4.5,1) {Ekuivalen: $\neg q \to \neg p$};
  \node[box=green!10] (c3) at (4.5,0) {Asumsikan $\neg q$, turunkan $\neg p$};
  \draw[thick, ->, >=stealth] (c1) -- (c2);
  \draw[thick, ->, >=stealth] (c2) -- (c3);
  \node[below=0.15cm of c3] {\textbf{Kontraposisi}};
\end{tikzpicture}
\caption{Alur bukti langsung (kiri) vs bukti kontraposisi (kanan)}
\label{fig:direct-vs-contrapositive}
\end{figure}

Bukti langsung paling efektif ketika anteseden $p$ memberikan informasi atau struktur yang memudahkan penurunan $q$ \cite{rosen2019}. Contoh: untuk membuktikan ``Jika $n$ genap maka $n^2$ genap'', asumsikan $n = 2k$ untuk suatu bilangan bulat $k$, lalu $n^2 = 4k^2 = 2(2k^2)$ yang jelas genap. Di sini asumsi $n$ genap memberi bentuk eksplisit $n = 2k$ yang langsung memungkinkan perhitungan. Sebaliknya, jika $p$ tidak memberikan ``pegangan'' yang memadai untuk menurunkan $q$, mungkin lebih baik mempertimbangkan metode lain seperti bukti kontraposisi \cite{forallx,stanford_intrologic}.

\textbf{Bukti kontraposisi} memanfaatkan ekuivalensi logika $p \rightarrow q \equiv \neg q \rightarrow \neg p$ \cite{rosen2019,copi2014}. Alih-alih membuktikan ``jika $p$ maka $q$'', kita membuktikan ``jika $\neg q$ maka $\neg p$''. Keduanya ekuivalen, sehingga membuktikan yang satu cukup untuk menetapkan yang lain. Kontraposisi sering dipilih ketika $\neg q$ lebih mudah ditangani daripada $p$; misalnya, ``$n$ ganjil'' ($\neg q$) dapat ditulis sebagai $n = 2k+1$, sementara ``$n^2$ genap'' ($p$) mungkin kurang memberi struktur untuk dikerjakan secara langsung \cite{levin_discrete,libretexts_proofs}.

Strategi pemilihan metode bergantung pada bentuk pernyataan dan kemudahan manipulasi \cite{rosen2019,levin_discrete}. Jika $q$ atau kondisi yang menyertainya lebih ``konkret'' (misalnya bentuk aljabar eksplisit) saat kita asumsikan $p$, gunakan bukti langsung. Jika $\neg q$ memberi bentuk yang lebih mudah (misalnya ``ganjil'' $\equiv$ ``bentuk $2k+1$''), gunakan kontraposisi. Tidak ada aturan baku; pengalaman dan eksplorasi awal membantu memilih pendekatan yang efisien \cite{forallx,copi2014}. Tabel~\ref{tab:direct-vs-contrapositive} merangkum perbandingan kedua metode \cite{libretexts_proofs,rosen2019}.

\begin{table}[htbp]
\centering
\small
\begin{tabular}{|>{\raggedright\arraybackslash}p{2.2cm}|>{\raggedright\arraybackslash}p{4cm}|>{\raggedright\arraybackslash}p{4cm}|}
\hline
\textbf{Aspek} & \textbf{Bukti langsung} & \textbf{Bukti kontraposisi} \\
\hline
Bentuk yang dibuktikan & $p \to q$: asumsikan $p$, tunjukkan $q$ & Ekuivalen $\neg q \to \neg p$: asumsikan $\neg q$, tunjukkan $\neg p$ \\
\hline
Kapan dipakai & $p$ memberi struktur/pegangan (mis.\ $n$ genap $\Rightarrow n=2k$) & $\neg q$ lebih mudah (mis.\ $n$ ganjil $\Rightarrow n=2k+1$) \\
\hline
Contoh pernyataan & ``Jika $n$ genap maka $n^2$ genap'' & ``Jika $n^2$ genap maka $n$ genap'' \\
\hline
\end{tabular}
\caption{Perbandingan bukti langsung dan kontraposisi}
\label{tab:direct-vs-contrapositive}
\end{table}

Dalam informatika, bukti langsung dan kontraposisi digunakan untuk verifikasi program dan spesifikasi formal \cite{stanford_cs103,mit_proof_methods}. Precondition dan postcondition metode program sering berbentuk implikasi: ``jika input memenuhi P maka output memenuhi Q''. Membuktikan kebenaran spesifikasi tersebut menjamin bahwa program berperilaku benar untuk input valid. Alat verifikasi formal dan proof assistants secara implisit mengandalkan teknik-teknik pembuktian ini ketika memeriksa kebenaran kode \cite{mit_6042,berkeley_logic}.

\begin{contoh}
Buktikan: ``Jika $n^2$ genap maka $n$ genap.'' Dengan kontraposisi, kita buktikan: ``Jika $n$ ganjil maka $n^2$ ganjil.'' Misalkan $n$ ganjil, maka $n = 2k+1$ untuk suatu bilangan bulat $k$. Maka $n^2 = (2k+1)^2 = 4k^2 + 4k + 1 = 2(2k^2+2k) + 1$, yang merupakan bilangan ganjil. Jadi jika $n$ ganjil maka $n^2$ ganjil; karena kontraposisi ekuivalen dengan implikasi asli, terbukti bahwa jika $n^2$ genap maka $n$ genap.
\end{contoh}

\begin{contoh}
Buktikan dengan bukti langsung: ``Jika $3n+2$ ganjil maka $n$ ganjil.'' Asumsikan $3n+2$ ganjil, maka $3n+2 = 2m+1$ untuk suatu bilangan bulat $m$. Dari sini $3n = 2m - 1$, sehingga $3n$ ganjil. Karena perkalian ganjil dengan genap menghasilkan genap, dan $3$ ganjil, $n$ harus ganjil agar $3n$ ganjil. Jadi $n$ ganjil. Terbukti.
\end{contoh}

\begin{contoh}
\textbf{Verifikasi numerik ``jika $n$ genap maka $n^2$ genap''} \cite{python_docs_boolean,levin_discrete}

Berikut verifikasi untuk beberapa nilai $n$ (bukan pengganti bukti formal, tetapi ilustrasi bahwa implikasi konsisten). Fungsi \code{genap} dan pengecekan postcondition menggambarkan ide bukti langsung: asumsikan $n$ genap, maka $n^2$ genap.

\begin{lstlisting}[style=pythonstyle, caption={Verifikasi numerik implikasi ``n genap $\Rightarrow$ n² genap''}]
def genap(n):
    return n % 2 == 0

# Postcondition: jika n genap maka n*n genap
for n in [2, 4, 6, 10, 0]:
    assert genap(n)  # asumsikan n genap
    assert genap(n * n)  # maka n^2 genap
print("Verifikasi: untuk n genap yang diuji, n^2 genap.")

# Kontraposisi: jika n^2 ganjil maka n ganjil
for n in [1, 3, 5, 7]:
    assert not genap(n * n)  # n^2 ganjil
    assert not genap(n)      # n ganjil
print("Kontraposisi: untuk n^2 ganjil, n ganjil.")
\end{lstlisting}
\end{contoh}

\begin{konsep}
Bukti langsung: asumsikan $p$, tunjukkan $q$ dengan langkah logis. Bukti kontraposisi: buktikan $\neg q \rightarrow \neg p$ sebagai gantinya dari $p \rightarrow q$. Pilih bukti langsung jika $q$ mudah diturunkan dari $p$; pilih kontraposisi jika $\neg q$ memberi bentuk yang lebih mudah untuk dikerjakan.
\end{konsep}
